\documentclass[11pt,twoside]{article}
% Shared formatting for all notes in this repository.
% Adapted from the UC Berkeley EECS CS267 lecture-note template:
% https://www.cs.cmu.edu/~ggordon/10725-F12/template.tex
% Personal-note adaptation: boxed title/author/date; ordinary numbering;
% existing natbib citations and mathematical environments are retained.
% Load after \documentclass; compile from the repository root.
\usepackage[T1]{fontenc}
\usepackage[utf8]{inputenc}
\usepackage{amsmath,amsfonts,amssymb,mathrsfs,amsthm}
\usepackage{graphicx,wrapfig}
\usepackage[inline,shortlabels]{enumitem}
\usepackage{centernot}
\usepackage[numbers]{natbib}
\usepackage{adjustbox,booktabs,arydshln}
\usepackage{caption,subcaption}
\usepackage{xcolor,tikz,pifont}
\usetikzlibrary{decorations.pathreplacing,angles,quotes,shapes.geometric,
  arrows.meta,positioning,calc,fit,backgrounds,shadows}
\usepackage[ruled,vlined]{algorithm2e}
\usepackage{xr}
\usepackage{hyperref}
\hypersetup{colorlinks=true,linkcolor=blue,citecolor=blue,urlcolor=blue}

% Berkeley EECS lecture-note page layout and paragraph spacing.
\setlength{\oddsidemargin}{0.25in}
\setlength{\evensidemargin}{-0.25in}
\setlength{\topmargin}{-0.6in}
\setlength{\textwidth}{6.5in}
\setlength{\textheight}{8.5in}
\setlength{\headsep}{0.75in}
\setlength{\parindent}{0in}
\setlength{\parskip}{0.1in}

% A wrapping box accommodates long research-note titles.
\makeatletter
\renewcommand{\maketitle}{%
  \pagestyle{myheadings}%
  \thispagestyle{plain}%
  \begin{center}
    \fbox{%
      \begin{minipage}{\dimexpr\textwidth-2\fboxsep-2\fboxrule\relax}
        \vspace{2mm}
        \textbf{Research Notes}\hfill\@date\par
        \vspace{2mm}
        {\centering\Large\bfseries\@title\par}
        \vspace{2mm}
        {\itshape Author: \@author\par}
        \vspace{2mm}
      \end{minipage}%
    }
  \end{center}
  \markboth{Research Notes}{Research Notes}%
  {\footnotesize\itshape
    Style adapted from the UC Berkeley EECS lecture-note template.\par}
  \vspace{4mm}%
}
\makeatother

\newcommand{\bR}{{\boldsymbol R}}
\newcommand{\bX}{{\boldsymbol X}}
\newcommand{\bY}{{\boldsymbol Y}}
\newcommand{\by}{{\boldsymbol y}}
\newcommand{\bx}{{\boldsymbol x}} 
\newcommand{\bz}{{\boldsymbol z}}
\newcommand{\bZ}{{\boldsymbol Z}}
\newcommand{\bH}{{\boldsymbol H}}
\newcommand{\bb}{{\boldsymbol b}}
\newcommand{\be}{{\boldsymbol e}}
\newcommand{\bbf}{{\boldsymbol f}}
\newcommand{\bg}{{\boldsymbol g}}
\newcommand{\bm}{{\boldsymbol m}}
\newcommand{\bw}{{\boldsymbol w}}
\newcommand{\bW}{{\boldsymbol W}}
\newcommand{\bv}{{\boldsymbol v}}
\newcommand{\bt}{{\boldsymbol t}}
\newcommand{\bA}{{\boldsymbol A}}
\newcommand{\ba}{{\boldsymbol a}}
\newcommand{\bD}{{\boldsymbol D}}
\newcommand{\bE}{{\boldsymbol E}}
\newcommand{\bi}{{\boldsymbol i}}
\newcommand{\bc}{{\boldsymbol c}}
\newcommand{\bC}{{\boldsymbol C}}
\newcommand{\bk}{{\boldsymbol k}}
\newcommand{\bK}{{\boldsymbol K}}
\newcommand{\bP}{{\boldsymbol P}}
\newcommand{\mK}{\mathcal{K}}
\newcommand{\bV}{{\boldsymbol V}}
\newcommand{\bG}{{\boldsymbol G}}
\newcommand{\mG}{\mathcal{G}}
\newcommand{\mH}{\mathcal{H}}
\newcommand{\mE}{\mathbb{E}}
\newcommand{\mX}{\mathcal{X}}
\newcommand{\mM}{\mathcal{M}}
\newcommand{\bU}{{\boldsymbol U}}
\newcommand{\bu}{{\boldsymbol u}}
\newcommand{\bI}{{\boldsymbol I}}
\newcommand{\bJ}{{\boldsymbol J}}
\newcommand{\bs}{{\boldsymbol s}}

\newcommand{\bxi}{{\boldsymbol \xi}}
\newcommand{\balpha}{{\boldsymbol \alpha}}
\newcommand{\bbeta}{{\boldsymbol \beta}}
\newcommand{\bgamma}{{\boldsymbol \gamma}}
\newcommand{\btheta}{{\boldsymbol \theta}}
\newcommand{\bTheta}{{\boldsymbol \Theta}}
\newcommand{\bvartheta}{{\boldsymbol \vartheta}}
\newcommand{\bepsilon}{{\boldsymbol \epsilon}}
\newcommand{\bvarepsilon}{{\boldsymbol \varepsilon}}
\newcommand{\bomega}{{\boldsymbol \omega}}
\newcommand{\bmu}{{\boldsymbol \mu}}
\newcommand{\bSigma}{{\boldsymbol \Sigma}}

\newcommand{\corr}{\mbox{Corr}}
\newcommand{\var}{{\mbox{Var}}}
\newcommand{\Var}{{\mbox{Var}}}
\newcommand{\cov}{{\mbox{Cov}}}
\newcommand{\diag}{{\mbox{diag}}}
\newcommand{\bUmis}{{\boldsymbol U}_{{\rm mis}}}
\newcommand{\bYmis}{{\boldsymbol Y}_{{\rm mis}}}
\newcommand{\bYobs}{{\boldsymbol Y}_{{\rm obs}}}
\newcommand{\tbXmis}{{\tilde{\boldsymbol X}}_{{\rm  mis}}}
\newcommand{\tbxmis}{{\tilde{\boldsymbol x}}_{{\rm  mis}}}
\newcommand{\bXmis}{{\boldsymbol X}_{{\rm mis}}}
\newcommand{\bXobs}{{\boldsymbol X}_{{\rm obs}}}
\newcommand{\bZmis}{{\boldsymbol Z}_{{\rm mis}}}
\newcommand{\bZobs}{{\boldsymbol Z}_{{\rm obs}}}
\newcommand{\tXmis}{{\tilde{X}}_{{\rm  mis}}}
\newcommand{\txmis}{{\tilde{x}}_{{\rm  mis}}}
\newcommand{\Xmis}{X_{{\rm  mis}}}
\newcommand{\Xobs}{X_{{\rm  obs}}}
\newcommand{\bymis}{{\boldsymbol y}_{{\rm  mis}}}
\newcommand{\byobs}{{\boldsymbol y}_{{\rm  obs}}}
\newcommand{\ymis}{y_{{\rm  mis}}}
\newcommand{\yobs}{y_{{\rm  obs}}}

\newcommand{\JRSSB}{Journal of the Royal Statistical Society: Series B}
\newcommand{\indep}{\rotatebox[origin=c]{90}{$\models$}}
\newcommand{\red}{\textcolor{red}}
\newcommand{\blue}{\textcolor{blue}}
\newcommand{\notetodo}[1]{\textcolor{blue}{[TODO: #1]}}

% Preserve the existing independent, section-based theorem counters.
\newtheorem{theorem}{Theorem}[section]
\newtheorem{lemma}{Lemma}[section]
\newtheorem{corollary}{Corollary}[section]
\newtheorem{formula}{Formula}[section]
\newtheorem{proposition}{Proposition}[section]
\theoremstyle{definition}
\newtheorem{definition}{Definition}[section]
\newtheorem{example}{Example}[section]
\newtheorem{assumption}{Assumption}
\theoremstyle{remark}
\newtheorem{remark}{Remark}

% Journal abbreviations used by the existing bibliography.
\def\AJE{American Journal of Epidemiology}
\def\ANNALS{Annals of Statistics}
\def\BIOK{Biometrika}
\def\BMCS{Biometrics}
\def\STATMED{Statistics in Medicine}
\def\STATSCI{Statistical Science}
\def\JSPI{Journal of Statistical Planning \& Inference}
\def\JRSSB{Journal of the Royal Statistical Society, Series B}
\def\COMMS{Communications in Statistics, Series A}
\def\JQT{Journal of Quality Technology}
\def\TECH{Technometrics}
\def\JASA{Journal of the American Statistical Association}
\def\CSDA{Computational Statistics \& Data Analysis}
\def\JCGS{Journal of Computational and Graphical Statistics}
\def\JCB{Journal of Computational Biology}
\def\BIOINF{Bioinformatics}
\def\JAMA{Journal of the American Medical Association}
\def\JNUTR{Journal of Nutrition}
\def\LETTERS{Letters in Probability and Statistics}
\def\JABES{Journal of Agricultural, Biological and Environmental Statistics}
\def\SCAN{Scandinavian Journal of Statistics}
\def\PNAS{Proceedings of the National Academy of Sciences USA}
\def\JMLR{Journal of Machine Learning Research}
\def\JSCS{Journal of Statistical Computation and Simulation}


\title{Continuous Time Markov Chains}
\author{Wen-Hung Wang}
\date{\today}

\begin{document}
\maketitle

\section{Introduction}
In this note, we review some properties of the Continuous Time Markov Chains (CTMCs). We first review some definitions and basic properties of CTMCs in \emph{discrete state space}. The content of this part is mostly from a STAT 532 note. 

After that we proceed to continuous state space. The purpose is to better understand the proof of convergence of Langevin diffusion to the Gibbs distribution in $W_2$. We take the content from Appendix A in \cite{raginsky2017non}. Note that, \cite{raginsky2017non} proves a \emph{nonasymptotic bound} for the $W_2$ convergence of SGLD with a constant step size.

\section{Discrete State Space}
A continuous time stochastic process $X$ consists of a sequence of random variables $X(t), t \geq 0$ evolving in a prescribed manner. Moreover, $X$ is right-continuous. We say $X$ is in a finite state space if $X(t) \in S$ where $S$ is countable. 

\begin{definition}[Markov]
    $X$ is Markovian if 
    \[
    P(X(t_n) = j | X(t_1)=i_1,...,X(t_{n-1})=i_{n-1}) = P(X(t_n) = j | X(t_{n-1})=i_{n-1})
    \]
    for all $0 \leq t_1 <...<t_n < \infty$ and $i_1,...,i_{n-1},j \in S$.
\end{definition}

The difficulty of CTMCs compared to DTMCs is that there is no time unit, meaning no exact analogy of transition matrix $P$: $P$ is the transition for a time unit; however, in CTMCs, $P_t$ depends on the time length $t$, which defines a family $\{P_t, t\geq 0\}$ called a transition semigroup, which will be defined later.

There are two approaches for studying CTMCs, one is to use infinitesimal calculus (e.g, generator) while the other is to study the embedded chain (e.g., chain at inter-arrival times). Examples of Markov CTMCs include Poisson and birth processes. 

\begin{definition}[Transition Probabilities and Time Homogeneity] Let 
\[
p_{ij}(t) = P(X(t) = j | X(0) = i), \quad t \geq 0
\]
be the transition probabilities of $X$. Then we say $X$ is time-homogeneous if $p_{ij}(s, t) = P(X(t) = j | X(s) = i) = p_{ij}(t-s)$. That is, the transition depends only on the length of time $t-s$ but not the initial time $s$.
\end{definition}

We assume that $X$ is time-homogeneous thereafter. 

\begin{definition}[Transition Semigroup and Stochastic Semigroup]
The family $\{P_t, t \geq 0\}$ is called a transition semigroup where $P_t = (p_{i,j}(t))_{i,j \in S}$ for any $t$. $\{P_t, t \geq 0\}$ is a stochastic semigroup if
\begin{enumerate}
    \item $P_0 = Id$.
    \item $p_{ij}(t) \geq 0$ and $\sum_{j \in S} p_{ij}(t)  = 1$ for any $t$.
\end{enumerate}
 Furthermore, $P_t$ satisfies the Chapman-Kolmogorov identity (semigroup property)
\[
P_{s+t} = P_s P_t
\]
for all $s, t \geq 0$.
\end{definition}

There at any given time $t$, $P_t$ is a stochastic matrix. Note that $P_t$ varies with $t$; therefore it is different from the transition matrix in DTMCs. However, we wish to have some "transition matrix" similar to the unit time transition $P$. The first approach is the infinitesimal calculus method, i.e., studying $P_h$ by considering $h \approx 0$.

\begin{definition}[Standard Semigroup]
    $P_t$ is called standard if it is continuous at $t=0$, i.e.,
    \[
    \lim_{t \to 0}P_t = P_0 = Id.
    \]
    Furthermore, by the CK identity,
    \[
    \lim_{h \to 0}P_{t+h} = \lim_{h \to 0} P_t P_h = P_t,
    \]
    i.e., $P_t$ is continuous at any time $t$.
\end{definition}

\subsection{The Generator and Its Properties}
We assume $P_t$ is standard thereafter. Under regularity conditions, consider a first-order Taylor expansion on $P_h$ around $0$, elementwisely,
\begin{align*}
    p_{ij}(h) &= p_{ij}(0) + q_{ij}h + {o(h)} \\
    &= \delta_{ij} + q_{ij}h + {o(h)}
\end{align*}
where $\delta_{ij} = 1$ if $i=j$; 0 otherwise. 

The meaning of $q_{ij}$ can be understood in the following way: In a small time $h$, the chain jumps from $i$ to $j$ with probability $q_{ij}h + o(h)$ while stays at $i$ with probability $1+q_{ii}h + o(h)$. Note that $q_{ij}$ relates to the rates in Poisson and birth processes.

\begin{definition}[Generator]
    A generator $Q$ of $\{P_t, t \geq 0\}$ is defined as $Q = (q_{ij})_{i, j \in S}$.
\end{definition}

Since $p_{ij}(h)$ is stochastic, we have $q_{ij} \geq 0$ if $i \neq j$ and $q_{ii} \leq 0$, i.e., the diagonal of $Q$ is non-positive. One can view $Q$ as a continuous time analogy of $P$ in discrete time; but crucially, it is not a stochastic matrix itself. 

Nevertheless, $Q$ provides an important way to define a CTMC, which we will discuss later. We summarize several properties of $Q$ as follows.

\begin{theorem}[Properties of $Q$, infinitesimal calculus] Here are some properties of $Q$:
    \begin{enumerate}
        \item \textbf{Differential.} $Q$ is the differential of $P_t$ at time $0$:
        \[
        \lim_{h \to 0}\frac{P_h-P_0}{h} = \lim_{h \to 0}\frac{P_h-Id}{h} = P_t'|_{t=0} = Q.
        \]
        \item \textbf{Orthogonality wrt one vector.} Since $\sum_{j \in S} p_{ij}(h) = 1$ for any $h \geq 0$ and $i \in S$, $\sum_{j \in S} q_{ij}=0$ for all $i \in S$. In matrix form, $Q\boldsymbol{1} = \boldsymbol{0}$.
        \item \textbf{Forward Equations.} 
        \[
        P_t'=P_tQ
        \]
        \item \textbf{Backward Equations.} 
        \[
        P_t'=QP_t
        \]
        \item \textbf{Transition from Generator.}
        From the forward and backward equations, we can express $P_t$ as a matrix exponential
        \[
        P_t = e^{tQ} = \sum_{k=0}^\infty \frac{t^kQ^k}{k!}.
        \]
    \end{enumerate}
    \textcolor{red}{One may use eigen-decomposition to solve for $P_t$ by solving the subsystem.}
\end{theorem}
\begin{proof}
    We prove properties 1, 3, and 5. From the Taylor expansion of $p_{ij}(h)$, we have 
    \[
    p_{ij}(h) = \delta_{ij} + q_{ij}h + {o(h)}.
    \]
    Rearranging terms, we have 
    \[
    \frac{p_{ij}(h)  - \delta_{ij}}{h} = q_{ij} + {o(1)}.
    \]
    Taking the limit $h \to 0$ proves 1.

    The forward equation can be written elementwisely as 
    \[
    p'_{ij}(t) = \sum_{k \in S}p_{ik}(t) q_{kj},
    \]
    which is our goal.
    By the CK identity, $P_{t+h} = P_t P_h$, i.e., 
    \begin{align*}
        p_{ij}(t+h) &= \sum_{k \in S}p_{ik}(t)p_{kj}(h) \\
        & = \sum_{k \in S}p_{ik}(t)[\delta_{kj} + q_{kj}h + {o(h)}] \\
        &= \sum_{k \neq j}p_{ik}(t)q_{kj}h  + p_{ij}(h) +  p_{ij}(h)q_{jj}h + {o(h)} \\
        &= p_{ij}(h) + \sum_{k \in S}p_{ik}(t)q_{kj}h + o(h).
    \end{align*}
    Rearranging terms, we obtain
    \[
    \frac{p_{ij}(t+h)-p_{ij}(h)}{h} = \sum_{k \in S}p_{ik}(t)q_{kj} + o(1).
    \]
    Taking limits $h \to 0$ on both sides, we conclude the forward equations. Backward equations can be proved similarly.

    To prove 5, we leverage the backward equations. Consider
    \begin{align*}  
        P_t - P_s &= \int_s^t P_u'du \\
        &= Q\int_s^t (P_u\textcolor{blue}{-P_{s}+ P_{s}})du \\
        &= QP_{s}(t-s) + Q\int_s^t P_u-P_{s}du \\
        &= QP_{s}(t-s) + Q\int_s^t \int_s^{u} P'_{u_2}d{u_2}du \\
        &= QP_{s}(t-s) + Q^2\int_s^t \int_s^{u} (P_{u_2} -P_{s}+P_{s})d{u_2}du \\
        &= QP_{s}(t-s) + Q^2P_{s}\underbrace{\int_s^t \int_s^{u} d{u_2}du}_{\frac{(t-s)^2}{2}}+ Q^2\int_s^t \int_s^{u} (P_{u_2} -P_{s})d{u_2}du \\
        &= QP_{s}(t-s) + Q^2P_{s}\frac{(t-s)^2}{2}+ Q^2\int_s^t \int_s^{u} \int_s^{u_2} P'_{u_3}d{u_3}d{u_2}du \\
        &= QP_{s}(t-s) + Q^2P_{s}\frac{(t-s)^2}{2}+ Q^3\int_s^t \int_s^{u} \int_s^{u_2} (P_{u_3}-P_s+P_s)d{u_3}d{u_2}du \\
        &= QP_{s}(t-s) + Q^2P_{s}\frac{(t-s)^2}{2}+ Q^3P_s\underbrace{\int_s^t \int_s^{u} \int_s^{u_2} d{u_3}d{u_2}du}_{\frac{(t-s)^3}{2\times 3}} + \int_s^{u_2} (P_{u_3}-P_s)d{u_3}d{u_2}du \\
        &\vdots
    \end{align*}
    Keeping iterating and setting $s=0$ gives the result. The key is to systematically use the forward equations to transform $P_t'$ to $QP_t$ and write the difference as an integral.
\end{proof}

Next, we discuss CTMCs in an embedded chain manner. We first define the inter-arrival time, given $X(s)=i$
\[
U=\inf_{t\geq 0}\{X(s+t) \neq i\}.
\]
Then we find the distribution of $U$ given $X(s)=i$.

\begin{theorem}[Inter-Arrival Time is Exponential]
    \[
    P(U>t | X(s)=i) = e^{q_{ii}t},
    \]
    i.e., $U  | X(s)=i \sim Exp(-q_{ii})$.
\end{theorem}

\begin{proof}
    We begin by showing $U$ is exponentially distributed. Recall that a continuous random variable is memory-less if and only if it is exponential, $P(U>s+t | U>s) = P(U>t)$. 
    \begin{align*}
        P(U>s+t | U>s, X(s)=i) &= P(U>s+t | X(s+s)=i, X(s)=i) \\
        &= P(U>s+t | X(s+s)=i) \\
        &= P(U>t | X(s)=i)
    \end{align*}
    The second equality follows from the Markov property while the last one follows from time homogeneity.

    Then we proceed to figure out the rate. In a small time interval $[t,t+h)$, 
    $P(X \text{ jumps}) = \lambda h +o(h)$. Here,
    \[
    P(X \text{ jumps}) = \underbrace{\sum_{k \neq i}q_{ij}}_{=\lambda}h + o(h).
    \]
    By Property 2, 
    \[
    \lambda = \sum_{k \neq i}q_{ij} = - q_{ii}.
    \]
\end{proof}

In what follows, we define an embedded Markov chain and its transition. 

\begin{definition}[Embedded Markov Chains]
    An embedded Markov chain is a discrete-time Markov chain constructed from a continuous-time process by observing the system only when it changes states.
\end{definition}

\begin{theorem}
For $i \neq j$,
    \[
    P(X \text{ jumps to } j |X \text{ jumps}, X(0) = i) = -\frac{q_{ij}}{q_{ii}}.
    \]
    Moreover, $\sum_{i \neq j} -\frac{q_{ij}}{q_{ii}} = 1$, i.e., $(P)_{ij} = \frac{q_{ij}}{q_{ii}}$ if $i \neq j$; 0 is $i=j$.
\end{theorem}

\begin{proof}
    In a small time interval with length $h$, by definition of the generator, one has
    \begin{align*}
        P(X \text{ jumps to } j |X \text{ jumps}, X(0) = i) &= \frac{P(X \text{ jumps to } j | X(0) = i)}{P(X \text{ jumps} | X(0) = i)} \\
        &= \frac{q_{ij}h+o(h)}{\sum_{j \neq i}q_{ij}h+o(h)} \\
        & \approx  \frac{q_{ij}}{-q_{ii}}. 
     \end{align*}
\end{proof}

\begin{remark}
    The generator $Q$ is the differential of $P_h$ in a small time interval at zero and can be seen as an analogy of the transition matrix in DTMCs but not a stochastic matrix itself. Given $Q$, the transition of a CTMC can be determined by $P_t=e^{tQ}$. Given $P_t$, taking derivative at 0 gives $Q$. Finally, the diagonal elements of $Q$ relate to the rate of inter-arrival time and the transition matrix of an embedded Markov chain.   
\end{remark}

\subsection{Classification of Chains}
We define and discuss irreducibility, stationary distribution, and ergodic theorem (limit theorem) of CTMCs with discrete state space.

\begin{definition}[Irreducibility]
    $X$ is irreducible if $p_{ij}(t) \geq 0$ for some $t \geq 0$. One can draw a graph of $Q$ to determine if $X$ is irreducible (connected between any two nodes?).
\end{definition}

\begin{definition}[Stationary Distribution]
    Let $\pi$ be a row vector. $\pi$ is a stationary distribution of $X$ if
    \begin{enumerate}
        \item $\pi_j \geq 0$ for all $j$ and $\sum_{j \in S}\pi_j=1$.
        \item For all $t \geq 0$,
        \[
        \pi = \pi P_t,
        \]
        i.e., $\pi_j = \sum_{i \in S}\pi_i p_{ij}(t)$ for any $t \geq 0$ and all $j \in S$.
    \end{enumerate}
\end{definition}

Compared to DTMCs, CTMCs require the definition to hold true for any $t$. The relation between $Q$ and $\pi$ is summarized below.

\begin{theorem} Let $Q$ be the generator of $X$.
    \[
    \pi \text{ is an invariant distribution of } X\Leftrightarrow \pi Q= 0
    \]
\end{theorem}

\begin{proof}
    \begin{align*}
        \pi = \pi P_t &\Leftrightarrow \pi = \pi e^{tQ} \\
        &\Leftrightarrow \pi = \pi \sum_{k=0}^\infty \frac{t^kQ^k}{k!} = \pi  + \pi \sum_{k=1}^\infty \frac{t^kQ^k}{k!}\\
        & \Leftrightarrow \pi Q =0
    \end{align*}
\end{proof}

Therefore, given $Q$, one can obtain $\pi$ by solving $\pi Q = 0$. The next theorem states some basic properties of $\pi$.

\begin{theorem}[Ergodic Theorem]
    Let $X$ be irreducible. 
    \begin{enumerate}
        \item If there exists stationary distribution $\pi$, then $\pi$ is unique and 
        \[
        \lim_{t \to \infty} p_{ij}(t) = \pi_j
        \]
        for all $i,j \in S$. That is, the chain forgets its initial value and converges to $\pi$ as $t \to \infty$.

        \item If there is no $\pi$, then
        \[
        \lim_{t \to \infty} p_{ij}(t) = 0
        \]
        for all $i,j \in S$.
    \end{enumerate}
\end{theorem}

\section{Continuous State Space}
\clearpage
\bibliography{ref}
\bibliographystyle{abbrv}
\end{document}
